\documentclass[10pt,a4paper]{amsart}
\usepackage[margin=19mm]{geometry}
\usepackage{amsmath,amssymb,mathtools}
\usepackage[scr=rsfs,cal=euler]{mathalfa}
\usepackage{xfrac}
\usepackage[unicode,hidelinks,bookmarks=false]{hyperref}
\newcommand{\ie}{i.e.\ }
\newcommand{\cf}{cf.\ }
\newcommand{\resp}{respectively\ }
\newcommand{\Subsection}{\subsection}
\providecommand{\nobreakdash}{\nobreak\hbox{-}}
\setlength{\emergencystretch}{2em}
\multlinegap0pt
\usepackage{bm}
\usepackage{bbm}
\usepackage{dsfont}
\usepackage{xspace}
\usepackage{cancel}
\usepackage{mathtools}
\usepackage{amsthm}
\makeatletter
\@ifundefined{theorem}{\newtheorem{theorem}{Theorem}}{}
\@ifundefined{lemma}{\newtheorem{lemma}[theorem]{Lemma}}{}
\makeatother
\def\N{\mathbb N}
\def\cU{\mathcal U}
\def\card{\operatorname{card}}
\def\dens{\operatorname{dens}}
\def\dist{\operatorname{dist}}
\def\meas{\operatorname{meas}}
\newcommand{\diam}{\operatorname{diam}}
\title[Hausdorff dimension in quasiregular dynamics]{Hausdorff dimension in quasiregular dynamics}
\author{Walter Bergweiler, Athanasios Tsantaris}
\date{Source: arXiv:2205.03626v1 (CC BY 4.0)}
\begin{document}
\maketitle

\noindent\emph{Source and licence.} Hausdorff dimension in quasiregular dynamics. Version arXiv:2205.03626v1 (CC BY 4.0), distributed under the Creative Commons Attribution 4.0 International licence, as recorded in the arXiv licence metadata for this exact version and shown on the abstract page https://arxiv.org/abs/2205.03626v1. Full rights notice: licenses/arxiv-2205.03626-CC-BY-4.0.txt.\par\smallskip

\noindent\textit{Editorial scope.} Counted unit: the Lemma whose source label is \texttt{lemma3} in the article (source0.tex lines 1318--1347, source label \texttt{lemma3}), delivered together with the article's own complete original proof of it (source0.tex lines 1348--1475, one \texttt{proof} environment of 128 lines). No mathematics is written, repaired or re-derived: statement and proof are the article's own text line for line. Required context retained uncounted: def-beta (lines 528--530); 7a0 (lines 551--553); 7a (lines 555--560); 7b (lines 562--565); 7a1 (lines 568--572); 5a (lines 1233--1235); 5a1 (lines 1237--1239). Every definition, display and label of those lines is retained unchanged. Omitted from this package and not redistributed: 1--527, 531--550, 554--554, 561--561, 566--567, 573--1232, 1236--1236, 1240--1317, 1476--2020. Nothing inside a retained excerpt is omitted or altered: every retained body is delivered exactly as the article's lines are written, so no omission receipt of any kind accompanies this package. Citations: the retained excerpts carry no citation command at all, so no bibliography entry of the article is transcribed and no reference is redistributed. Reference adaptation: none was needed and none was made; every reference inside the delivered unit resolves inside it, so no unresolved reference or citation is produced. Printed slips kept literal: no printed word, symbol, hypothesis or number of the counted statement or of its proof was altered. Layout and class: the article's own class is \texttt{amsart} with packages including inputenc, amsmath, amsfonts, hyperref, array, booktabs, amssymb, mathtools; the wrapper is the lane's pinned \texttt{amsart} editorial wrapper at 10pt on A4, which restates the theorem-like environments the retained text needs and adds the article's own macro definitions that the retained text needs (\texttt{\textbackslash N}, \texttt{\textbackslash cU}, \texttt{\textbackslash card}, \texttt{\textbackslash dens}, \texttt{\textbackslash diam}, \texttt{\textbackslash dist}, \texttt{\textbackslash meas}), with extra wrapper packages (bm, bbm, dsfont, xspace, cancel, mathtools). The delivered numbering is the unit's own. Assets: no retained span contains a figure environment, a graphics-inclusion command, a PDF-page inclusion, a table or a TikZ picture; no image file is copied into this package, the package declares no asset dependency and no image file is redistributed with it. Licence: version arXiv:2205.03626v1 of this article is distributed under the Creative Commons Attribution 4.0 International licence, as recorded in the arXiv licence metadata for this exact version and shown on the abstract page https://arxiv.org/abs/2205.03626v1. A full rights notice is delivered at licenses/arxiv-2205.03626-CC-BY-4.0.txt. Characters: every retained excerpt is pure ASCII, so the delivered engine, XeLaTeX with its default Latin Modern text font, reports no missing character.
\begin{equation}\label{def-beta}
\beta=\frac{3-\rho}{2}
\end{equation}

\begin{equation}\label{7a0}
\card(S\cap B(x,r)) \leq \alpha_0 r^{2-2\beta}.
\end{equation}

\begin{equation}\label{7a}
\frac{1}{64} |x|^{-2\beta}r^2
\leq \card(S\cap B(x,r))
\leq 6 |x|^{-2\beta}r^2
\quad\text{for}\ 8|x|^\beta\leq r\leq \frac12 |x|
\end{equation}

\begin{equation}\label{7b}
\card(S\cap B(x,r))\leq 324
\quad\text{for}\ r<8|x|^\beta.
\end{equation}

\begin{equation}\label{7a1}
\card(S\cap B(x,r))
\leq 324 |x|^{-2\beta}r^2
\quad\text{for}\ |x|^\beta\leq r\leq 8|x|^\beta  .
\end{equation}

\begin{equation}\label{5a}
\diam P\leq C_2
\end{equation}

\begin{equation}\label{5a1}
\meas P\geq C_3.
\end{equation}

\begin{lemma} \label{lemma3}
There exists $C_4,C_5,C_6,R>0$ such that if $F\in\{G,H\}$ and $t\geq R$, then
\begin{equation}\label{3h1}
\card \cU_F(t)
\geq C_4 \frac{t^{\rho}}{\eta(t)}
\end{equation}
and
\begin{equation}\label{3h2}
\dens\!\left(U_F(t),K(t)\right)
\geq C_5 \frac{t^{\rho-3}}{\eta(t)}.
\end{equation}
If  $x\in Q(t)$, then
\begin{equation}\label{3h3}
\card\{P\in  \cU_F(t)\colon P\cap B(x,\tau)\neq\emptyset\}
\leq
\begin{cases}
C_6(\tau +1) &\text{if}\ \tau\leq t^{\beta},\\
C_6\tau^\rho    &\text{if}\ \tau> t^{\beta},\\
\end{cases}
\end{equation}
with $\beta$ defined by~\eqref{def-beta}.
In particular 
\begin{equation}\label{eq3h3}
	\card \cU_F(t)\leq C_6(2t)^\rho.
\end{equation}
Moreover, if $P,P'\in \cU_F(t)$, then
\begin{equation}\label{3h4}
\dist(P,P')\geq\frac14 \eta(t).
\end{equation}
\end{lemma}

\begin{proof}
Recall that  $B_2((x_1,x_2),r)$ denotes
the planar disk centred at $(x_1,x_2)$ and of radius $r$. Assuming that $R$ is large, it follows from~\eqref{5a} that 
if 
\begin{equation}\label{5e}
s
\in
 B_2\!\left( \frac15(t, t),\frac{t}{31}\right)
\subset
 B_2\!\left( \frac15(t, t),\frac{t}{30}-C_2\right)
\end{equation}
and 
\begin{equation}\label{5f}
\frac{3t}{4}+C_2\leq k\eta(k) \leq \frac{5t}{4}-C_2 ,
\end{equation}
then $P_F(s,k)\subset Q(t)$ and thus $P_F(s,k)\in \cU_F(t)$.
By~\eqref{7a}, 
the cardinality of the set of all $s\in S$ satisfying~\eqref{5e} 
is at least 
\begin{equation}\label{5g}
\frac{1}{64} \left(\frac{\sqrt{2}t}{5}\right)^{-2\beta} \frac{t^2}{31^2}
=\frac{2^{-\beta} }{64\cdot 5^{-2\beta} \cdot 31^2} \cdot t^{2-2\beta},
\end{equation}
provided $t$ is sufficiently large.

To estimate the cardinality of the set of all $k$ satisfying~\eqref{5f} we 
claim that~\eqref{5f} is satisfied if 
\begin{equation}\label{5g1}
\left(\frac{3t}{4}+C_2\right)\frac{1}{\eta(t)-1}
\leq k\leq
\left(\frac{5t}{4}-C_2\right)\frac{1}{\eta(t)}
\end{equation}
To see this note that~\eqref{5g1} implies that 
$\sqrt{t}\leq k\leq t$ and thus $\eta(t)-1\leq\eta(k)\leq \eta(t)$.
Using~\eqref{5a} we can now deduce~\eqref{5f} from~\eqref{5g1}.

The cardinality of the 
set of all $k\in\N$ which satisfy~\eqref{5g1} and hence also~\eqref{5f} is at least
\begin{equation}\label{5h}
\frac13 \frac{t}{\eta(t)},
\end{equation}
provided $t$ is sufficiently large.
Combining the bounds~\eqref{5g} and~\eqref{5h} we find that 
\begin{equation}\label{5i}
\card \cU_F(t)
\geq \frac{2^{-\beta} }{3\cdot 64\cdot 5^{-2\beta} \cdot 31^2} \frac{t^{3-2\beta}}{\eta(t)}
\end{equation}
for large $t$. 
Now~\eqref{3h1} follows from~\eqref{def-beta}.

Combining~\eqref{3h1} with~\eqref{5a1} we obtain~\eqref{3h2} with $C_5=C_3C_4/(6\pi)$.

To prove~\eqref{3h3} let $x\in Q(t)$. Arguments similar to the ones used
to obtain~\eqref{5e} and~\eqref{5f} show that if $P_F(s,k)\in \cU_F(t)$ with 
$P_F(s,k)\cap B(x,\tau)\neq\emptyset$, then
\begin{equation}\label{5j}
s\in 
B_2((x_1,x_2), \tau+C_2)
\end{equation}
and 
\begin{equation}\label{5k}
|k\eta(k)-x_3|\leq \tau+C_2.
\end{equation}
By the definition of~$Q(t)$ we have
\begin{equation}\label{5l}
\frac{t}{6}\leq |(x_1,x_2)|\leq |x| \leq 2t .
\end{equation}
Suppose first that $t^\beta\leq\tau\leq t/24$. 
Then, provided $R$ is large enough,  $\tau+C_2\leq 2\tau$ and hence
$B_2((x_1,x_2), \tau+C_2)\subset B_2((x_1,x_2), 2\tau)$.
Since $2\tau\leq t/12\leq 2|(x_1,x_2)|$ we can use~\eqref{7a} and~\eqref{7a1}
to bound the cardinality of $S\cap B_2((x_1,x_2),2\tau)$.
This is bigger than 
the cardinality $N_1$ of the set of all $s\in S$ satisfying~\eqref{5j}. Thus
\begin{equation}\label{5m}
N_1\leq 324 |(x_1,x_2)|^{-2\beta} (2\tau)^2
\leq 324\cdot 6^{2\beta}\cdot  4 t^{-2\beta}\tau^2 .
\end{equation}
Suppose next that $\tau\geq t/24$. Then 
$N_1\leq \alpha_0 (2\tau)^{2-2\beta}\leq
\alpha_0 2^{2-2\beta} 24^{2\beta} t^{-2\beta}\tau^2$ by~\eqref{7a0}.
Suppose finally that $\tau\leq t^\beta$.
Then $\tau+C_2\leq t^\beta+C_2\leq 6^\beta|x|^\beta+C_2\leq 8|x|^\beta$ for large $R$
by~\eqref{5l}. Now~\eqref{7b} yields that $N_1\leq 324$. 
Altogether it follows that there exists $C_6>0$ such that 
\begin{equation}\label{5n2}
4 N_1\leq 
\begin{cases}
 C_6 &\text{if}\ \tau\leq t^{\beta},\\
\displaystyle C_6 t^{-2\beta}\tau^2   &\text{if}\ t^{\beta}\leq\tau .\\
\end{cases}
\end{equation}

For $s\in S$ we now consider the cardinality $N_2(s)$ of the set of all 
$k$ such that $P_F(s,k)$ intersects $B(x,\tau)$.
Since 
\begin{equation}\label{5m1}
(k+1)\eta(k+1)-k\eta(k)\geq \eta(k)\geq \eta(t)-1\geq \frac12 \eta(t)
\end{equation}
if $P_F(s,k)\in \cU_F(t)$ we deduce from~\eqref{5a} that
\begin{equation}\label{5n}
N_2(s)\leq N_2:=\frac{4(\tau+C_2)}{\eta(t)}  +1. 
\end{equation}
Combining the above estimates we obtain 
\begin{equation}\label{5n1}
\card\{P\in  \cU_F(t)\colon P\cap B(x,\tau)\neq\emptyset\}
\leq N_1 N_2.
\end{equation}

Suppose that $\tau\leq t^\beta$. 
Noting that~\eqref{5n} yields that $N_2\leq \tau+2$ we obtain~\eqref{3h3}
from~\eqref{5n1} and~\eqref{5n2}.

Suppose next that $\tau> t^\beta$.
Since $Q(t)\subset K(t)\subset B(x,2t)$ we may assume that $\tau\leq 2t$.
By~\eqref{5n} we have $N_2\leq \tau$ for large $t$.
Together with~\eqref{5n1} and~\eqref{5n2} this yields that 
\begin{equation}\label{5n3}
\card\{P\in  \cU_F(t)\colon P\cap B(x,\tau)\neq\emptyset\}
\leq \frac14 C_6 t^{-2\beta} \tau^3
=\frac14 C_6 \left(\frac{\tau}{t}\right)^{2\beta} \tau^\rho\leq C_6 \tau^\rho.
\end{equation}
So we obtain~\eqref{3h3} also in this case.
Since, as already mentioned, $Q(t)\subset B(x,2t)$, this also yields~\eqref{eq3h3}.

Finally, \eqref{3h4} follows from~\eqref{5a} and~\eqref{5m1}, provided
$R$ is large enough.
\end{proof}

\end{document}
